\CharacterSheet{enforcer}{Executora da lei}{

\shortdescription{A exigente guardiã da lei da família, cuja autoridade só é superada pela da chefe.}

\ObjectivesHeading{Objetivos}
\begin{itemize}
  \item Descobrir quem matou \Character{newcomer} e conseguir que essa pessoa seja expulsa por votação. (\PositiveReward{1}, \NegativeReward{1})
  \item Recuperar o livro de regras. (\PositiveReward{1}, \NegativeReward{1})
  \item Identificar casos em que membros da família tenham violado as leis do livro de regras. (\PositiveReward{1}, \NegativeReward{1}) por cada caso
\end{itemize}

\LosingConditionsHeading{Condições de derrota}
Não conseguir expulsar por votação quem matou \Character{newcomer}, caso haja uma pessoa responsável pelo homicídio.

\SecretsHeading{Segredos}
\begin{itemize}
  \item \Character{ghoul} foi transformado com o meu sangue e tem de obedecer a qualquer instrução que eu lhe dê como ordem.
  \item Ofereci algum tipo de imunidade em troca de alguma coisa, embora não me lembre dos pormenores.
\end{itemize}

\ActionsHeading{Ações}
\begin{itemize}
  \item Sê formal, fria e controladora. Transforma cada acusação num processo: exige nomes, horas, testemunhas e consequências.
  \item Quando precisares de alguma coisa de \Character{ghoul}, dá-lhe uma ordem clara. Ele não pode recusar, embora revelar esse poder possa causar problemas.
\end{itemize}

\BackgroundHeading{Contexto}

\timestamp{Há dois mil anos}, era uma humana normal com necessidades normais. Precisava de um mundo menos caótico, de pessoas com consideração pelos outros e de consequências para quem falhasse à sua comunidade.

Um dia, \Character{vampirehead} acolheu-me na sua família porque queria alguém que transformasse a sua vontade em lei. Essa continua a ser a noite e a festa mais especial da minha vida. Escreveu o \Clue{rulebook} ali mesmo: um membro que a família considere culpado de violar uma regra perde os poderes. A minha palavra teria mais peso do que a de qualquer outro, exceto a da chefe. A justiça e o controlo assentavam-me na perfeição.

Durante muito tempo, não tive grande relevância, já que \Character{vampirehead} não tinha pressa de acrescentar novos membros. Enquanto só existíamos \Character{vampirehead}, \Character{usurper} e eu, diverti-me com a imortalidade e os poderes, embora só pudesse fazê-lo à vontade longe da luz do sol.

\timestamp{Décadas depois da morte de Cristo}, \Character{buttler} foi acolhida numa festa semelhante à minha. Antes era comerciante; agora, criada de \Character{vampirehead}.

\timestamp{No ano 100 d.C.}, durante uma viagem, conheci \Character{ghoul}, um bobo humano encantador. Amava-o, mas as regras proibiam um vampiro de ter um amante humano. Pedi a \Character{vampirehead} que o transformasse. Pedi também que a poção dele fosse feita com o meu sangue. Disse à chefe que isso nos aproximaria. A verdade é que queria certezas: queria alguém que não pudesse escapar ao meu controlo. O contrato obriga \Character{ghoul} a obedecer a qualquer instrução que eu lhe dê como ordem.

Só \Character{ghoul} e eu assistimos à transformação. Depois, dei-lhe uma ordem permanente: teria de me contar, para sempre, todas as informações relevantes que descobrisse sobre qualquer outro membro da família. Tem sido uma excelente fonte de informações, quer goste do papel ou não. Transmite sempre a imagem de um artista desajeitado e engraçado, pelo que ninguém espera que trabalhe como espião.

Ao longo dos séculos, a família foi ganhando cada vez mais membros. Aprendi todas as regras, repetia-as nos encontros e escondi o livro em segurança no castelo. Confiavam em mim para esclarecer dúvidas, identificar infrações e impedir que a família se tornasse uma turba de imortais. Quando os novos membros se revelavam traiçoeiros, tratava das sentenças e retirava-lhes os poderes. Quando estava demasiado ocupada, \Character{vampirehead} encontrava outras soluções.

Desde então, passaram a entrar membros com menos frequência. Tinha mais tempo para passar com \Character{ghoul}. Gostávamos especialmente de criar música.

\Character{supremacist} era uma governante que foi acolhida para ajudar a escolher novos membros. Talvez porque, mesmo como humana, tratava bem a nossa família. Entrou em \timestamp{500 d.C.}.

\Character{lover} era um guerreiro que nos podia ajudar nas operações militares. Entrou em \timestamp{700 d.C.}.

\Character{prevlover} podia ajudar-nos nas investigações. Entrou em \timestamp{1000 d.C.}. Trabalhámos muitas vezes juntos. Sinceramente, era mais útil pelo artefacto que tinha do que pela inteligência. Embora já não me lembre bem de que artefacto era.

\timestamp{No ano 1400 d.C.}, entrou \Character{traitor}. Foi uma escolha controversa. Era uma espécie de benfeitora pública, o que não correspondia propriamente às necessidades da família. \Character{lover} insistiu até \Character{vampirehead} a aceitar. Eu não confiava nela, não confio e nunca confiarei.

\timestamp{Há dois anos}, \Character{vampirehead} apresentou um humano à família numa festa: \Character{newcomer}. Nada que merecesse atenção.

\timestamp{Há cerca de um ano}, alguém roubou uma página do livro de regras. Fiquei furiosa e pensei numa pessoa que parecia a suspeita perfeita. Por alguma razão, já não me lembro de quem era. Ainda assim, não queria precipitar-me. Pedi ajuda a \Character{prevlover} e seguimos ambos o rasto de mensagens incriminatórias na casa da pessoa suspeita. Não me lembro das mensagens. Era \Character{prevlover} quem tinha as nossas notas.

\timestamp{Há dois meses}, como era costume, \Character{vampirehead} marcou \timestamp{16/7/2026} como data da festa para acolher \Character{newcomer} na família. Todos os vampiros tinham de vir.

\timestamp{Há umas duas semanas}, \Character{vampirehead} perguntou se mais alguém, além dela, de \Character{buttler} e de \Character{newcomer}, queria ajudar. Disse que \Character{usurper} se tinha oferecido e chegaria ao castelo duas noites antes da festa.

\timestamp{Ontem, algum tempo depois das 16:00}, senti algo estranho. As regras que tinha apregoado e feito cumprir durante tanto tempo desapareceram-me da memória. Não consigo recordar a maior parte desse dia. Lembro-me vagamente de ter oferecido algum tipo de imunidade a um membro da família. Nem sequer me lembro de como escondi o livro de regras.

\GMNotes{
\paragraph{Porque se pode suspeitar desta personagem?}
\Character{enforcer} tem uma autoridade invulgar, escondeu o \Clue{rulebook} e controla secretamente \Character{ghoul} através de uma ordem mágica. A tarde de que não se lembra parece dar-lhe uma oportunidade para ocultar um motivo ou um ataque, enquanto a disponibilidade para aplicar castigos torna plausível a hipótese de um crime a sangue-frio.

\paragraph{Como se pode provar a inocência desta personagem?}
O incidente de perda de memória ocorreu às \timestamp{16:15}, horas depois de \Character{newcomer} ter recebido secretamente a primeira dose da poção de \Character{usurper} ao almoço. \Character{enforcer} não preparou o vinho nem participou no almoço. O testemunho de \Character{prevlover} e os resíduos do ritual da memória mostram que as ações esquecidas foram uma tentativa de impedir uma morte prevista, e não de a causar.

\paragraph{Qual é o papel desta personagem na história e no jogo?}
\Character{enforcer} dá uma voz processual à investigação. Pode transformar rumores em acusações, divulgar as consequências do \Clue{rulebook} quando for recuperado e dar uma ordem direta que obrigue \Character{ghoul} a revelar informações relevantes. A memória afetada também dá a \Character{prevlover} um crime grave para esconder, embora não seja um homicídio.

\paragraph{Como pode esta personagem encaminhar a investigação para a solução?}
Se os jogadores descobrirem o ritual da memória ou pressionarem \Character{prevlover}, estabelece que \Character{enforcer} estava a agir devido ao aviso de \Character{ghoul} sobre a morte de \Character{newcomer}. Uma ordem clara pode obrigar \Character{ghoul} a repetir o que contou. A cronologia recuperada afasta a atenção do ritual da memória, posterior, e leva-a para o almoço, em especial para o vinho fornecido por \Character{usurper}.
}
}
